\section{Results}
\label{sec:results}
Table~\ref{tab:results} collects every model that was evaluated. Cross-validated
scores (mean over the five folds) guided all choices; test scores were computed
once, for the baseline and the two final models only.

\begin{table}[ht]
  \centering
  \small
  \caption{Cross-validated (training set, 5 folds) and test-set scores. AUC:
  higher is better; log-loss: lower is better. Only the three final models were
  evaluated on the test set.}
  \label{tab:results}
  \begin{tabular}{@{}lcccc@{}}
    \toprule
    & \multicolumn{2}{c}{5-fold CV} & \multicolumn{2}{c}{Test set} \\
    \cmidrule(lr){2-3}\cmidrule(l){4-5}
    Model & AUC (mean $\pm$ sd) & Log-loss & AUC & Log-loss \\
    \midrule
    Baseline (constant 6.7\%)                & $0.500 \pm 0.000$ & 0.245 & 0.500 & 0.245 \\
    Logistic regression, raw features        & $0.851 \pm 0.004$ & 0.188 & --    & --    \\
    Logistic regression, $\log(1+x)$, ridge $C{=}10$ & $0.852 \pm 0.005$ & 0.185 & 0.858 & 0.183 \\
    Logistic regression, binned (experiment) & $0.855 \pm 0.005$ & 0.183 & --    & --    \\
    Random forest, min.\ leaf 20             & $0.863 \pm 0.003$ & 0.179 & 0.866 & 0.177 \\
    \bottomrule
  \end{tabular}
\end{table}

\paragraph{Logistic regression.} The log transform of the long-tailed features
improved the cross-validated AUC slightly (0.851 to 0.852) and the log-loss
more clearly (0.188 to 0.185). Figure~\ref{fig:regularization} shows the
penalty path. For $C \geq 0.1$, ridge and lasso give the same scores, and the
selected model ($C = 10$, ridge) is in practice unpenalized. Only strong
penalties change anything, and they hurt: at $C = 10^{-4}$ the lasso keeps at
most one of the 14 coefficients and falls to the baseline, while the ridge
still ranks clients well (AUC 0.847) but its log-loss, 0.286, is worse than
the constant baseline, because all probabilities are squeezed toward the same
value. A logistic regression on 20 quantile bins of the five continuous
features, an attempt to capture the non-linear curves of
Figure~\ref{fig:eda}, gained only 0.003 of AUC, less than one standard
deviation across folds, and was not kept.

\begin{figure}[ht]
  \centering
  \includegraphics[width=\linewidth]{fig_regularizacao}
  \caption{Logistic regression: cross-validated AUC (a) and log-loss (b)
  against the inverse penalty $C$ for ridge and lasso. The dotted line marks
  the selected $C$; the dashed line in (b) is the constant baseline.}
  \label{fig:regularization}
\end{figure}

\paragraph{Random forest.} Figure~\ref{fig:forest} shows the bias--variance
trade-off controlled by the minimum leaf size. With leaves of one client, the
trees memorize the training folds (training AUC 1.000) and the validation AUC
is the worst of the grid (0.843). Larger leaves make each tree simpler: the
training score falls and the validation score rises until about 20--50
clients per leaf, then starts to fall again. The selected leaf size (20,
lowest validation log-loss) gives a cross-validated AUC of 0.863.

\begin{figure}[ht]
  \centering
  \includegraphics[width=\linewidth]{fig_floresta}
  \caption{Random forest: AUC (a) and log-loss (b) on the training and
  validation folds as the minimum leaf size grows. The dotted line marks the
  selected value.}
  \label{fig:forest}
\end{figure}

\paragraph{Test set.} On the 30{,}000 test clients the logistic regression
reaches an AUC of 0.858 and the random forest 0.866
(Figure~\ref{fig:roc}), both within 0.007 of their cross-validated values. In
the Gini coefficient that credit-risk teams usually report ($2\,\mathrm{AUC}-1$),
this is 0.72 against 0.73. Over
most of the ROC curve the two models are almost indistinguishable; the forest
is ahead mainly between 10\% and 40\% of false positives. With the default
threshold of 0.5, both models beat the constant model's accuracy by less than
half a percentage point (93.7\% against 93.3\%), which is why accuracy was not
used.

\begin{figure}[ht]
  \centering
  \includegraphics[width=0.42\linewidth]{fig_roc}
  \caption{ROC curves of the two final models on the test set.}
  \label{fig:roc}
\end{figure}

\paragraph{Decision threshold.} The thresholds chosen on the out-of-fold
training predictions to reach a recall of 0.70 were 0.079 for the logistic
regression and 0.101 for the forest. Applied unchanged to the test set
(Table~\ref{tab:threshold}), both still catch about 70\% of defaulters, so the
choice generalized. For the same recall, the forest refuses 641 fewer good
payers (3{,}949 against 4{,}590 false positives, 14\% fewer), which is how its
higher AUC translates into a business decision. In the terms of a credit
desk, every client below the threshold is approved: the logistic regression
approves 80.0\% of the applicants (TN + FN) and the forest 82.1\%, and the bad
rate among approved clients (FN over approved) falls from 6.7\% in the whole
test portfolio to 2.4\% for both models. At the same bad rate, the forest
approves 649 more clients.

\begin{table}[ht]
  \centering
  \small
  \caption{Test-set confusion counts at the threshold chosen by
  cross-validation (recall $\geq 0.70$ on out-of-fold predictions). Positive =
  default.}
  \label{tab:threshold}
  \begin{tabular}{@{}lccccccc@{}}
    \toprule
    Model & Threshold & TP & FN & FP & TN & Precision & Recall \\
    \midrule
    Logistic regression & 0.079 & 1{,}421 & 584 & 4{,}590 & 23{,}405 & 0.236 & 0.709 \\
    Random forest       & 0.101 & 1{,}413 & 592 & 3{,}949 & 24{,}046 & 0.264 & 0.705 \\
    \bottomrule
  \end{tabular}
\end{table}

\paragraph{What drives the predictions.} Figure~\ref{fig:interpretation}
compares the logistic regression's odds ratios with the forest's impurity-based
importances. Both models rely on the same four features: the three
late-payment counts and revolving utilization. In the logistic regression, one
standard deviation more of (log) revolving utilization multiplies the odds of
default by 1.97, and one more standard deviation of age by 0.79. The flags
built during the exploratory analysis carry signal: the 96/98 code flag
multiplies the odds by 12.6 and the revolving-above-400\% flag by 2.8. The
forest ranks these flags low, since each one affects few clients.

\begin{figure}[ht]
  \centering
  \includegraphics[width=\linewidth]{fig_interpretacao}
  \caption{(a) Odds ratios $e^{\beta_j}$ of the logistic regression, per
  standard deviation for numerical features and for a flag switched on;
  values above 1 raise the risk. (b) Mean decrease in Gini impurity in the
  random forest (unsigned).}
  \label{fig:interpretation}
\end{figure}
